% =============== 导言区 ===============
\usepackage[UTF8]{ctex}             % 中文支持
% ===== 数学相关 =====
\usepackage{amsmath}                 % 数学公式环境（align、equation等）
\usepackage{amssymb}                 % 额外数学符号（如\mathbb、R等）
\usepackage{amsfonts}                % 数学字体（如\mathrm、\mathcal等）
\usepackage{mathtools}               % amsmath增强（箭头公式、矩阵环境等）
\usepackage{newtxtext}               % 英文正文字体（Times风格）
\usepackage{newtxmath}               % 英文数学字体（Times风格）

% ===== 页面与排版 =====
\usepackage{geometry}                % 页面尺寸和边距设置
\usepackage{setspace}                 % 行距控制

% ===== 表格与图片 =====
\usepackage{graphicx}                % 图片插入
\usepackage{xcolor}                   % 颜色支持（如 \textcolor{blue}{...}）
\usepackage{array}                   % 表格列格式控制
\usepackage{longtable}               % 跨页长表格
\usepackage{booktabs}                 % 专业三线表（\toprule、\midrule、\bottomrule）
\usepackage{multirow}                % 表格单元格跨行合并
\usepackage{wrapfig}                 % 文字环绕图片
\usepackage{textcomp}                % 特殊符号（如 \textcopyright、\textohm 等）

\usepackage{xparse}                % 带可选参数的自定义环境
\usepackage{environ}               % 环境体捕获（用于 answer 环境）
\usepackage{fancyhdr}              % 页眉页脚（习题集模板使用）
\usepackage{newunicodechar}        % 兼容少量数学题中常见的 Unicode 符号
\usepackage[colorlinks=false,pdfborder={0 0 0}]{hyperref} % 目录与跳转

\graphicspath{{./Images/}{Images/}{../Images/}} % 习题图片搜索路径

\newunicodechar{ⅰ}{\ensuremath{\mathrm{i}}}
\newunicodechar{ⅱ}{\ensuremath{\mathrm{ii}}}
\newunicodechar{ⅲ}{\ensuremath{\mathrm{iii}}}
\newunicodechar{ⅳ}{\ensuremath{\mathrm{iv}}}

% =============== 自定义命令 ===============
% 说明：
% 1. 命令用法中的 [...] 表示可选参数，可以省略；省略时使用默认值。
% 2. 命令用法中的 * 表示星号版本，例如 \QuestionSetSection*{...}，星号不写时使用自动编号版本。
% 3. 文件名、图片名等参数一般写在 {...} 中；图片标题等可选内容写在 [...] 中。

% ===== 题目版式命令 =====
% 选择题选项命令，可选参数控制排版方式；可选参数省略时默认为 [1]。
% 用法：\choices{A}{B}{C}{D} 等价于 \choices[1]{A}{B}{C}{D}，4列横排
% 用法：\choices[2]{A}{B}{C}{D}  2列2行
% 用法：\choices[4]{A}{B}{C}{D}  1列4行（每个选项一行）
\newcommand{\choices}[5][1]{%
  \par
  \ifcase#1\relax
  \or % 1: 4列横排
    \begin{tabular*}{0.8\linewidth}[t]{@{\extracolsep{\fill}}cccc}
      \hspace{2em}A．#2 & B．#3 & C．#4 & D．#5
    \end{tabular*}%
  \or % 2: 2列2行
    \begin{tabular}{p{0.45\linewidth}p{0.45\linewidth}}
      \hspace{2em}A．#2 & B．#3 \\
      \hspace{2em}C．#4 & D．#5
    \end{tabular}%
  \or % 3: 未使用
  \or % 4: 1列4行
    \begin{tabular}{l}
      \hspace{2em}A．#2 \\
      \hspace{2em}B．#3 \\
      \hspace{2em}C．#4 \\
      \hspace{2em}D．#5
    \end{tabular}%
  \fi
}

% 选择题选项括号命令
% 用法：\choiceBox 表示选项括号 (~~~~~~)，使用 \mbox 防止括号被换行拆开
\newcommand{\choiceBox}{\mbox{(~~~~~~)}}

% 下划线填空命令
% 用法：\blank{3cm} 表示长度为 3cm 的下划线，使用 \mbox 防止下划线被换行拆开
\newcommand{\blank}[1]{\underline{\mbox{\hspace{#1}}}}

% ===== 题内图片命令 =====
% 图片均为非浮动排版，适合放在 question 环境中的任意位置
\makeatletter
\newcommand{\MC@questionImageCaption}[1]{%
  \if\relax\detokenize{#1}\relax\else
    \par\smallskip
    {\small #1}%
  \fi
}

% 居中图片
% 用法：\questionImage{image1.png} 等价于 \questionImage[0.45\linewidth]{image1.png}
% 用法：\questionImage[0.6\linewidth]{image1.png} 指定图片宽度
% 用法：\questionImage[0.45\linewidth]{image1.png}[图1] 最后的 [图1] 为可选图片标题
\NewDocumentCommand{\questionImage}{O{0.45\linewidth} m o}{%
  \par\smallskip
  \begin{center}
    \includegraphics[width=#1]{#2}%
    \IfNoValueF{#3}{\MC@questionImageCaption{#3}}%
  \end{center}
  \smallskip
}

% 右图左文局部图文块
% 用法：\begin{questionImageRight}{image1.png}左侧文字\end{questionImageRight}
%      等价于 \begin{questionImageRight}[0.32\linewidth]{image1.png}左侧文字\end{questionImageRight}
% 用法：\begin{questionImageRight}[0.32\linewidth]{image1.png}[图1]左侧文字\end{questionImageRight}
%      [0.32\linewidth] 为右侧图片宽度，[图1] 为可选图片标题
\NewDocumentEnvironment{questionImageRight}{O{0.32\linewidth} m o +b}{%
  \par\smallskip
  \noindent
  \begin{minipage}[t]{\dimexpr\linewidth-#1-1em\relax}
    \vspace{0pt}\raggedright
    #4
  \end{minipage}%
  \hspace{1em}%
  \begin{minipage}[t]{#1}
    \vspace{0pt}\centering
    \includegraphics[width=\linewidth]{#2}%
    \IfNoValueF{#3}{\MC@questionImageCaption{#3}}%
  \end{minipage}%
  \par\smallskip
}{}

% 左图右文局部图文块
% 用法：\begin{questionImageLeft}{image1.png}右侧文字\end{questionImageLeft}
%      等价于 \begin{questionImageLeft}[0.32\linewidth]{image1.png}右侧文字\end{questionImageLeft}
% 用法：\begin{questionImageLeft}[0.32\linewidth]{image1.png}[图1]右侧文字\end{questionImageLeft}
%      [0.32\linewidth] 为左侧图片宽度，[图1] 为可选图片标题
\NewDocumentEnvironment{questionImageLeft}{O{0.32\linewidth} m o +b}{%
  \par\smallskip
  \noindent
  \begin{minipage}[t]{#1}
    \vspace{0pt}\centering
    \includegraphics[width=\linewidth]{#2}%
    \IfNoValueF{#3}{\MC@questionImageCaption{#3}}%
  \end{minipage}%
  \hspace{1em}%
  \begin{minipage}[t]{\dimexpr\linewidth-#1-1em\relax}
    \vspace{0pt}\raggedright
    #4
  \end{minipage}%
  \par\smallskip
}{}

% 双图并排
% 用法：\questionImagePair{imageA.png}{imageB.png}
%      等价于 \questionImagePair[0.45\linewidth]{imageA.png}{imageB.png}
% 用法：\questionImagePair[0.45\linewidth]{imageA.png}[图1]{imageB.png}[图2]
%      [0.45\linewidth] 为每张图片宽度，[图1]、[图2] 为可选图片标题
\NewDocumentCommand{\questionImagePair}{O{0.45\linewidth} m o m o}{%
  \par\smallskip
  \noindent
  \begin{minipage}[t]{#1}
    \vspace{0pt}\centering
    \includegraphics[width=\linewidth]{#2}%
    \IfNoValueF{#3}{\MC@questionImageCaption{#3}}%
  \end{minipage}%
  \hfill
  \begin{minipage}[t]{#1}
    \vspace{0pt}\centering
    \includegraphics[width=\linewidth]{#4}%
    \IfNoValueF{#5}{\MC@questionImageCaption{#5}}%
  \end{minipage}%
  \par\smallskip
}
\makeatother

% ===== 答案显示模式 =====
\newif\ifshowans
\newif\ifshowanalysis
\newif\ifchapteranswers
\newif\ifbookanswers
\showansfalse
\showanalysistrue
\chapteranswerstrue
\bookanswersfalse

% 开启题后显示答案功能（主要用于试卷入口 ExamPaper.tex）
% 用法：\showAnswers        显示答案和分析
% 用法：\showAnswersOnly    只显示答案，不显示分析
% 用法：\hideAnswers        隐藏题后答案
\newcommand{\showAnswers}{\global\showanstrue\global\showanalysistrue}
\newcommand{\showAnswersOnly}{\global\showanstrue\global\showanalysisfalse}
\newcommand{\hideAnswers}{\global\showansfalse}

% 习题集常用答案方案（主要用于 QuestionSet.tex，每次只选一种）
% 用法：\answersNone              不显示任何答案
% 用法：\answersByChapter         题后不显示答案，每章末尾显示答案速校
% 用法：\answersAtEnd             题后和章末不显示答案，全书末尾显示答案速校
% 用法：\answersInline            每题后显示答案和分析
% 用法：\answersInlineBrief       每题后只显示答案
% 用法：\answersInlineAndChapter  每题后显示答案和分析，并在章末显示答案速校
% 用法：\studentBook              等价于 \answersByChapter
% 用法：\teacherBook              等价于 \answersInline
\newcommand{\answersNone}{%
  \global\showansfalse
  \global\chapteranswersfalse
  \global\bookanswersfalse
}
\newcommand{\answersByChapter}{%
  \global\showansfalse
  \global\chapteranswerstrue
  \global\bookanswersfalse
}
\newcommand{\answersAtEnd}{%
  \global\showansfalse
  \global\chapteranswersfalse
  \global\bookanswerstrue
}
\newcommand{\answersInline}{%
  \global\showanstrue
  \global\showanalysistrue
  \global\chapteranswersfalse
  \global\bookanswersfalse
}
\newcommand{\answersInlineBrief}{%
  \global\showanstrue
  \global\showanalysisfalse
  \global\chapteranswersfalse
  \global\bookanswersfalse
}
\newcommand{\answersInlineAndChapter}{%
  \global\showanstrue
  \global\showanalysistrue
  \global\chapteranswerstrue
  \global\bookanswersfalse
}
\newcommand{\studentBook}{\answersByChapter}
\newcommand{\teacherBook}{\answersInline}

% 解答题空白区域命令
% 用法：\solutionSpace{5cm} 留出 5cm 答题空间；显示题后答案时自动不留大块空白
\newcommand{\solutionSpace}[1]{%
  \ifshowans
    \par\smallskip
  \else
    \vspace{#1}\par
  \fi
}

% 自动标号计数器
\newcounter{qno}
\newcounter{questionSetChapterNumber}
\newcounter{questionSetSectionNumber}

% 题目环境（替代 \Q{} 的推荐用法）
% 用法：\begin{question}题干内容\end{question} 自动递增题号
% 用法：\begin{question}[题源]题干内容\end{question} 可选参数 [题源] 会显示在题号后
% 使用 begin/end 环境，解决长题目花括号匹配困难的问题
\NewDocumentEnvironment{question}{O{}}{%
  \stepcounter{qno}%
  \raggedright
  \theqno．%
  \if\relax\detokenize{#1}\relax\else（#1）\fi
  \ignorespaces
}{\par}

% 重置题号命令
% 用法：\resetQuestionNumber 将题号重置为 0；通常由 \QuestionSetSection 自动调用
\newcommand{\resetQuestionNumber}{%
  \setcounter{qno}{0}%
}

% ===== 标题与答案命令 =====
\makeatletter
\newcommand{\chapteranswerlist}{}
\newcommand{\bookanswerlist}{}
\let\aclist\chapteranswerlist
\newif\ifsectionanswergroups
\sectionanswergroupsfalse

% 记录大题分组
\newcommand{\ACGroup}[1]{%
  \g@addto@macro\chapteranswerlist{%
    \par\smallskip
    \noindent\textbf{#1}\par
    \noindent
  }%
  \g@addto@macro\bookanswerlist{%
    \par\smallskip
    \noindent\textbf{#1}\par
    \noindent
  }%
}

% 记录单题答案
\newcommand{\ACItem}[2]{%
  \g@addto@macro\chapteranswerlist{%
    \textbf{#1.}~#2\qquad
  }%
  \g@addto@macro\bookanswerlist{%
    \textbf{#1.}~#2\qquad
  }%
}

% 大题标题命令
% 用法：\questionGroup{标题内容} 输出加粗大题标题，并把标题写入答案速校分组
\newcommand{\questionGroup}[1]{%
  \ifsectionanswergroups\else
    \ACGroup{#1}%
  \fi
  \noindent\textbf{#1}\par
}

% 答案环境（替代 \A{答案}{分析}）
% 用法：\begin{answer}{答案}分析内容\end{answer}
% 用法：\begin{answer}{答案}\end{answer} 分析为空时不显示【分析】标签
% 第一个花括号参数会进入答案速校；环境正文是详细分析
\NewEnviron{answer}[1]{%
  \ifshowans
    \par\smallskip
    {\color{red}\textbf{【答案】}}{\boldmath\textbf{#1}}\par
    \ifshowanalysis
      \if\relax\detokenize\expandafter{\BODY}\relax\else
        {\color{red}\textbf{【分析】}}~\BODY\par
      \fi
    \fi
    \vspace{6ex}
  \fi
  \begingroup
    \edef\x{\endgroup
      \noexpand\ACItem{\theqno}{\unexpanded{#1}}%
    }%
  \x
}

% 参考答案页 / 本章答案速校
% 用法：\printAnswerCheck 显示当前已收集的答案速校内容
\newcommand{\printAnswerCheck}{%
    \begingroup
    \small
    \noindent{\large \textbf{答案速校}}\par
    \ifx\chapteranswerlist\@empty
      \noindent（本部分暂无答案。）\par
    \else
      \noindent\chapteranswerlist\par
    \fi
    \endgroup
    \medskip
}

% 用法：\printBookAnswerCheck 在全书末尾输出全书答案速校；仅 \answersAtEnd 模式下生效
\newcommand{\printBookAnswerCheck}{%
  \ifbookanswers
    \clearpage
    \begin{center}
      {\Large\bfseries 全书答案速校}
    \end{center}
    \begingroup
    \small
    \ifx\bookanswerlist\@empty
      \noindent（本书暂无答案。）\par
    \else
      \noindent\bookanswerlist\par
    \fi
    \endgroup
  \fi
}
\makeatother

% ===== 章节结构命令 =====

% 习题集封面、目录与页眉
% 用法：\UseQuestionSetStyle{高中数学习题集} 设置页眉、页脚和页码样式
\newcommand{\UseQuestionSetStyle}[1]{%
  \setlength{\headheight}{15pt}%
  \pagestyle{fancy}
  \fancyhf{}
  \fancyhead[L]{#1}
  \fancyhead[R]{\nouppercase{\leftmark}}
  \fancyfoot[C]{\thepage}
  \renewcommand{\headrulewidth}{0.4pt}
}

% 用法：\QuestionSetCover{主标题}{副标题}{底部信息} 生成习题集封面
\newcommand{\QuestionSetCover}[3]{%
  \begin{titlepage}
    \centering
    \vspace*{0.16\textheight}
    {\Huge\bfseries #1\par}
    \vspace{1.2cm}
    {\Large #2\par}
    \vfill
    {\large #3\par}
  \end{titlepage}
}

% 用法：\QuestionSetTOC 输出目录并自动分页
\newcommand{\QuestionSetTOC}{%
  \clearpage
  \tableofcontents
  \clearpage
}

% 用法：\includeChapter{01-集合与常用逻辑用语}
%      自动输入 chapters/01-集合与常用逻辑用语.tex，不需要写 chapters/ 和 .tex
\newcommand{\includeChapter}[1]{%
  \input{chapters/#1.tex}%
}

\makeatletter
\newcommand{\MC@appendBookChapter}[1]{%
  \g@addto@macro\bookanswerlist{%
    \par\bigskip
    \noindent{\large\bfseries #1}\par
    \smallskip
  }%
}

% 章标题（新页开始，居中大号）
% 用法：\QuestionSetChapter{函数与基本性质} 自动显示“第1章　函数与基本性质”
% 用法：\QuestionSetChapter*{附录} 显示“附录”，不推进章节编号
% 非星号版本会重置小节编号、写入目录、更新页眉，并开始收集本章答案
\NewDocumentCommand{\QuestionSetChapter}{s m}{%
  \clearpage
  \gdef\chapteranswerlist{}%
  \global\sectionanswergroupsfalse
  \IfBooleanTF{#1}{%
    \def\MC@chapterTitle{#2}%
  }{%
    \stepcounter{questionSetChapterNumber}%
    \setcounter{questionSetSectionNumber}{0}%
    \def\MC@chapterTitle{第\arabic{questionSetChapterNumber}章　#2}%
  }%
  \expandafter\MC@appendBookChapter\expandafter{\MC@chapterTitle}%
  \phantomsection
  \addcontentsline{toc}{section}{\MC@chapterTitle}%
  \markboth{\MC@chapterTitle}{\MC@chapterTitle}%
  \begin{center}
    {\Large\bfseries \MC@chapterTitle}
  \end{center}
}

% 节标题（知识点名称，左对齐加粗）
% 用法：\QuestionSetSection{函数的概念与表示} 自动显示“1.1　函数的概念与表示”
% 用法：\QuestionSetSection*{本章综合训练} 显示“本章综合训练”，不推进小节编号
% 非星号版本会递增小节编号；两种版本都会重置题号、写入目录并开启答案分组
\NewDocumentCommand{\QuestionSetSection}{s m}{%
  \par\medskip
  \resetQuestionNumber
  \global\sectionanswergroupstrue
  \IfBooleanTF{#1}{%
    \def\MC@sectionTitle{#2}%
  }{%
    \stepcounter{questionSetSectionNumber}%
    \def\MC@sectionTitle{\arabic{questionSetChapterNumber}.\arabic{questionSetSectionNumber}　#2}%
  }%
  \expandafter\ACGroup\expandafter{\MC@sectionTitle}%
  \phantomsection
  \addcontentsline{toc}{subsection}{\MC@sectionTitle}%
  \noindent{\large\bfseries \MC@sectionTitle}\par
  \smallskip
}

% 本章答案速校 — 显示累积答案后清空，各章答案互不干扰
% 用法：\printChapterAnswers 输出本章答案速校；仅章末答案模式开启时显示
\newcommand{\printChapterAnswers}{%
  \ifchapteranswers
    \clearpage
    \begin{center}
      {\Large\bfseries 本章答案速校}
    \end{center}
    \printAnswerCheck
  \fi
  \gdef\aclist{}%
  \gdef\chapteranswerlist{}%
}
\makeatother
